\section{Related Work}
\label{sec:related}

\paragraph{Minimum-edge-cut redistricting.}
The algorithmic redistricting literature has converged on edge-cut minimisation
as a mathematically clean, politically neutral criterion
\citep{chikina2017,deford2019,duchin2020}.
METIS multilevel graph partitioning \citep{karypis1998} is the standard solver.
\citet{b7paper} establish that for states where the seed-sweep converges, the
MEC plan is a property of the graph, not the seed.
GeoSection \citep{b8paper} introduces isoperimetric normalisation to select
among feasible split ratios at each bisection level.

\paragraph{Compactness and geographic fairness.}
Polsby-Popper \citep{polsby1991}, Reock \citep{reock1961}, and convex-hull ratio
are standard single-district compactness metrics used in court challenges.
These metrics evaluate the shape of individual districts, not the structure of
the bisection tree.
Edge-cut minimisation is related but distinct: it minimises total boundary
length across all districts simultaneously, which is a stronger and
more globally consistent criterion.

Area balance as a redistricting criterion has received limited formal attention.
\citet{vickrey1961} proposes that legislative districts should be ``compact and
contiguous'' without specifying land area formally.
\citet{niemi1990} study compactness empirically but focus on shape metrics.
To our knowledge, AreaSection is the first algorithm to enforce simultaneous
population and land-area balance as a formal METIS dual constraint.

\paragraph{Multi-constraint graph partitioning.}
METIS supports $k$-way partitioning with $\mathtt{ncon} > 1$ balance constraints
\citep{karypis1999metis}.
The dual-constraint mode has been used in scientific computing for
load balancing with memory constraints \citep{devine2002,hendrickson1995},
but has not previously been applied to redistricting.
We use $\mathtt{ncon}=2$ with $\mathtt{tpwgts}$ specifying population fractions
and $\mathtt{ubvec}=[1.001,\, 1.10]$ specifying per-constraint imbalance
tolerances (0.1\,\% for population, 10\,\% for area).

\paragraph{Urban concentration and partisan effects.}
The geographic sorting of American voters --- Democrats concentrated in urban
cores, Republicans spread through suburban and rural areas --- is well
documented \citep{rodden2019,chen2013}.
\citet{stephanopoulos2015} propose the efficiency gap as a measure of partisan
bias arising from geographic sorting.
\citet{mcghee2014} show that geographic sorting alone, without any intentional
gerrymandering, produces Republican-leaning maps under most compactness
criteria.
AreaSection is designed to test whether enforcing land-area balance
mitigates this systematic effect.
Our empirical results (\S\ref{sec:evaluation}) show it changes the
competitiveness structure but not the seat allocation, suggesting that
geographic sorting is robust to changes in the bisection ratio structure.

\paragraph{Lorenz curves in spatial analysis.}
The Lorenz curve \citep{lorenz1905} measures concentration in distributions.
\citet{frank2010} apply spatial Lorenz curves to population-area distributions
in urban analysis.
We adapt this tool to redistricting feasibility analysis: the population-area
Lorenz curve (population fraction as a function of cumulative area fraction,
sorted densest-first) predicts which split ratios can simultaneously satisfy
population and area constraints.
